\chapter{AI trong quản lý và phân tích dữ liệu}

\subsection{Tóm tắt}

Chương 4 khám phá cách AI đang cách mạng hóa việc quản lý và phân tích dữ liệu học thuật bằng cách tự động hóa các tác vụ tốn nhiều công sức như xử lý và chuẩn hóa dữ liệu (data wrangling), làm sạch, tích hợp và phân tích khám phá. Từ dữ liệu khảo sát có cấu trúc đến nội dung văn bản phi cấu trúc, các công cụ AI tạo sinh và học máy có thể phát hiện bất thường, đề xuất cách hài hòa hóa dữ liệu, xây dựng các đặc trưng mới và làm lộ ra các mẫu hình ở quy mô lớn. Chương cung cấp những hiểu biết thực tiễn về các công cụ như Google DataPrep, Airtable và Wikidata, đồng thời nhấn mạnh nhu cầu tiếp tục có sự giám sát của con người. Chương cảnh báo về việc quá phụ thuộc vào các gợi ý tự động và nhấn mạnh tầm quan trọng của việc kiểm chứng, khả năng tái lập và quản trị dữ liệu có đạo đức trong các quy trình nghiên cứu có sự hỗ trợ của AI.

\subsection{Từ khóa}

Làm sạch dữ liệu, AI trong nghiên cứu, tích hợp dữ liệu, phân tích khám phá, kỹ thuật đặc trưng (feature engineering)

Nghiên cứu học thuật đương đại đã được định hình lại một cách căn bản bởi quy mô và độ phức tạp chưa từng có của dữ liệu hiện có (Kitchin, 2014; Boyd and Crawford, 2012). Hầu như mọi lĩnh vực học thuật đều đang bị dữ liệu hóa (Flensburg and Lomborg, 2021): dữ liệu hóa (datafication) là từ thời thượng mới, và dữ liệu trở thành đồng tiền mới của giới học thuật.

Sau khi các nhà nghiên cứu có được bộ dữ liệu của mình, thông qua thu thập truyền thống, các phương pháp tính toán hoặc hệ thống tự động, họ phải đối mặt với những thách thức đáng kể trong việc chuẩn bị, tích hợp và tiến hành các phân tích sơ bộ đối với dữ liệu. Những giai đoạn chuẩn bị này đòi hỏi chuyên môn và nguồn lực đáng kể, đồng thời có thể tiềm ẩn các lỗ hổng do sai sót hệ thống, thiên lệch hoặc sự thiếu hiệu quả trong quy trình nghiên cứu (Kaisler et al., 2013; Jagadish et al., 2014).

Để ứng phó với những thách thức này, các công nghệ trí tuệ nhân tạo, đặc biệt là học máy và xử lý ngôn ngữ tự nhiên (NLP), đã trở thành công cụ quan trọng trong quản lý dữ liệu nghiên cứu. Các hệ thống này hiện hỗ trợ cấu trúc hóa bộ dữ liệu, phát hiện bất thường ở quy mô lớn và chỉ ra những hướng phân tích đầy triển vọng (Jordan and Mitchell, 2015; LeCun, Bengio, and Hinton, 2015). Nhờ giảm thời gian dành cho việc chuẩn bị dữ liệu thường quy, nhà nghiên cứu có thể dành nhiều sự chú ý hơn cho phân tích chuyên sâu và diễn giải lý thuyết, qua đó nâng cao hiệu quả và chất lượng công việc của mình.

Lợi ích của việc đưa trí tuệ nhân tạo vào nghiên cứu vượt ra ngoài việc cải thiện hiệu quả. Các hệ thống được huấn luyện trên những bộ dữ liệu lớn rất xuất sắc trong việc phân loại thông tin phi cấu trúc, phát hiện các mẫu hình ẩn và nhận diện những điểm bất nhất tinh vi mà các phương pháp truyền thống có thể bỏ sót (Chen, Chiang, and Storey, 2012; Wu et al., 2014). Năng lực này đặc biệt có giá trị trong nghiên cứu học thuật hiện đại, nơi các bộ dữ liệu trải rộng từ tài liệu lịch sử đến các dòng thông tin số theo thời gian thực.

Việc triển khai các công nghệ này trong nghiên cứu đòi hỏi sự cân nhắc và hoạch định cẩn trọng. Nhà nghiên cứu phải bảo đảm sự phù hợp giữa đặc điểm dữ liệu, cách tiếp cận phương pháp luận và các giải pháp công nghệ được lựa chọn. Cần đặc biệt chú ý duy trì tính minh bạch, khả năng tái lập và quản trị dữ liệu có đạo đức (Kitchin, 2014; Mittelstadt and Floridi, 2016). Khi được triển khai đúng cách, các phương pháp quản lý dữ liệu tiên tiến này có thể cải thiện đáng kể chất lượng nghiên cứu, đồng thời tạo điều kiện cho quá trình chuyển từ thông tin thô sang hiểu biết khoa học có ý nghĩa. Chương này sẽ tìm hiểu cách khai thác những công cụ đó để phục vụ lợi ích của chúng ta.

\section{Xử lý và làm sạch dữ liệu với sự hỗ trợ của AI}

Xử lý và chuẩn hóa dữ liệu (data wrangling) cùng với làm sạch dữ liệu là một trong những giai đoạn tốn nhiều nguồn lực nhất của nghiên cứu (tiện thể, cũng là một trong những giai đoạn nhàm chán nhất!), đòi hỏi học giả phải giải quyết nhiều vấn đề khác nhau, từ định dạng tệp không nhất quán đến các mục nhập thiếu hoặc sai sót (Kandel et al., 2011). Ngay cả những sơ suất nhỏ cũng có thể làm suy yếu giá trị của các phân tích tiếp theo (Zhu and Wu, 2014). Với những bộ dữ liệu ngày nay trải rộng trên bảng tính, cơ sở dữ liệu và các nguồn cấp dữ liệu mạng xã hội, việc chuẩn hóa nhãn thủ công hoặc rà soát hàng nghìn dòng để tìm điểm bất thường có thể trở nên quá tải. Các nhà nghiên cứu ngày càng sử dụng các kỹ thuật dựa trên AI để tiết kiệm thời gian và phát hiện, sửa lỗi đáng tin cậy hơn.

Nhiều công cụ mã nguồn mở và thương mại cho thấy AI có thể tự động hóa hoặc bán tự động hóa các tác vụ như loại bỏ bản ghi trùng lặp, điền giá trị thiếu và phát hiện bất thường (Hellerstein et al., 2017). Chẳng hạn, Trifacta Wrangler (hiện đã được đổi thương hiệu thành Alteryx Designer Cloud), vốn khởi nguồn từ nghiên cứu của Kandel và các cộng sự, áp dụng học máy để đề xuất các phép biến đổi dữ liệu, trong khi OpenRefine có thể gợi ý những phép gộp tiềm năng cho các giá trị chuỗi không nhất quán (ví dụ: ``U.S.A.'' so với ``United States'') bằng cách phân tích mức độ tương đồng giữa chúng. Trong Python, các gói như pandas-profiling hay great-expectations tạo ra các báo cáo làm nổi bật những vấn đề tiềm ẩn về chất lượng dữ liệu, giúp nhà nghiên cứu nhanh chóng phát hiện các phân phối bất thường hoặc cột bị gán nhãn sai trước khi những sai sót đó làm méo mó phân tích của họ.

Một lợi ích then chốt của các công cụ có sự hỗ trợ của AI là khả năng học từ phản hồi của người dùng và tinh chỉnh các đề xuất của chúng. Giả sử bạn tải một bảng tính lớn lên Trifacta Wrangler, hoặc chạy một tập lệnh làm sạch trong OpenRefine, hoặc dùng Google DataPrep, rồi tự tay điều chỉnh một số trường văn bản, chẳng hạn thống nhất ``Not applicable'', ``N/A'' và các ô trống về một mã duy nhất. Khi bạn xác nhận những thay đổi này, công cụ có thể tự động áp dụng chúng trên toàn bộ tập dữ liệu, giúp giảm đáng kể việc thiết lập quy tắc lặp đi lặp lại (Ebraker and Hellerstein, 2005). Quy trình thích ứng này đặc biệt hữu ích trong các môi trường dữ liệu trực tiếp, nơi những lần làm mới định kỳ có thể đưa vào các điểm bất thường mới cần được xử lý nhất quán.

Hãy xét một tình huống trong đó bạn gộp các tập dữ liệu từ nhiều trường đại học để so sánh tỷ lệ tốt nghiệp. Một cơ sở gắn nhãn hồ sơ chưa hoàn tất là ``WD'', một cơ sở khác là ``Withdrawn'', và một cơ sở nữa là ``Drop''. Một nền tảng có AI có thể gom nhóm các thuật ngữ này dựa trên độ tương đồng ngữ nghĩa, rồi nhắc bạn chấp nhận hoặc đổi tên chúng đồng loạt thành ``Withdrawn''. Việc chấp thuận lời nhắc đó sẽ kích hoạt sự đồng nhất tự động trên mọi trường liên quan. Tương tự, trong một nghiên cứu về khối lượng công việc của giảng viên, bạn có thể dùng pandas-profiling trong Python để phát hiện rằng cột số ``Hours\_Taught'' của một trường đại học chứa các giá trị ngoại lai, chẳng hạn con số đáng ngờ ``2,000 hours'', cho thấy có lỗi nhập liệu hoặc gán nhãn đơn vị. Việc đánh dấu sớm những điểm bất thường này giúp nhà nghiên cứu sửa chúng trước khi chúng ảnh hưởng đến việc xây dựng mô hình về sau.

Ngoài dữ liệu có cấu trúc, các giải pháp dựa trên AI còn hỗ trợ việc làm sạch dữ liệu phi cấu trúc như bản ghi phỏng vấn hoặc câu trả lời khảo sát mở (Klinkenberg 2017). Các công cụ như xử lý ngôn ngữ tự nhiên (NLP) dựa trên văn bản trong Python (ví dụ spaCy hoặc NLTK) có thể loại bỏ văn bản rập khuôn, các tuyên bố miễn trừ trách nhiệm dư thừa hoặc dấu thời gian, cho phép nhà nghiên cứu tập trung vào các lời kể của người tham gia. Trong các dự án định tính quy mô lớn, bạn có thể đưa các tài liệu PDF vào một quy trình do AI vận hành để tự động lọc bỏ các tiêu đề hành chính lặp lại hoặc các tuyên bố pháp lý không liên quan. Điều này bảo đảm bạn còn lại phần nội dung văn bản cốt lõi để mã hóa và phân tích chủ đề. Những quy trình này giúp việc chuẩn bị dữ liệu thông suốt hơn đồng thời giảm các lỗi của con người do phân tích cú pháp văn bản thủ công gây ra.

Suy cho cùng, việc tận dụng AI cho xử lý và làm sạch dữ liệu không chỉ đẩy nhanh quy trình nghiên cứu mà còn nâng cao chất lượng và độ tin cậy của các phát hiện về sau. Nhờ giảm thiểu các tác vụ lặp lại hoặc thủ công, nhà nghiên cứu có thể tập trung áp dụng chuyên môn của mình, như chất vấn các điểm bất thường, tinh chỉnh câu hỏi nghiên cứu và hình thành những diễn giải vững chắc. Tuy nhiên, mỗi phép biến đổi do AI thực hiện đều cần được ghi chép kỹ lưỡng, với nhật ký chi tiết về các đề xuất đã chấp nhận và các quy tắc ánh xạ, nhằm bảo đảm tính minh bạch và khả năng tái lập. Khi được cân bằng với sự giám sát của con người và việc kiểm chứng cẩn thận, làm sạch dữ liệu với sự hỗ trợ của AI mở đường cho việc khám phá dữ liệu tự tin hơn và xây dựng mô hình nâng cao, đặt nền tảng vững chắc cho những đóng góp khoa học có ý nghĩa.

\begin{examplebox}{Ví dụ thực tế: Làm sạch các phản hồi khảo sát khách hàng bằng Google DataPrep}

1. Thiết lập và nhập dữ liệu

Truy cập Google DataPrep

Trong Google Cloud Console, hãy điều hướng đến mục ``Dataprep'' (thường nằm trong ``Big Data'' hoặc ``Data Analytics'').

Nếu đây là lần đầu bạn sử dụng, hãy tạo một \textbf{luồng DataPrep (DataPrep flow)} hoặc dự án mới bằng cách làm theo hướng dẫn trên màn hình để kết nối với bucket Google Cloud Storage (GCS) của bạn.

Tải lên hoặc kết nối với bộ dữ liệu của bạn

Đặt các tệp thô của bạn (ví dụ: CSV, Excel hoặc JSON) vào một bucket Google Cloud Storage.

Trong DataPrep, nhấp vào \textbf{``Import Datasets''}, sau đó chọn nguồn của bạn. Công cụ sẽ tạo bản xem trước và tự động lấy mẫu dữ liệu.

Đánh giá dữ liệu ban đầu

DataPrep tự động tạo một \textbf{tổng quan về chất lượng dữ liệu}, làm nổi bật các vấn đề như giá trị bị thiếu hoặc không khớp kiểu dữ liệu.

Khám phá \textbf{hồ sơ (profile)} của từng cột (phân phối, giá trị nhỏ nhất/lớn nhất) để phát hiện các lỗi tiềm ẩn (ví dụ: ký hiệu bất thường, giá trị nằm ngoài phạm vi).

2. Chuẩn hóa văn bản

Xác định và chọn các cột văn bản tự do

Nhận diện các cột như ``Customer\_Comments'' hoặc ``Feedback\_Text.''

Nhấp vào tiêu đề cột để truy cập các tùy chọn biến đổi.

Sử dụng các hàm ``Text Cleanup'' tích hợp sẵn

\textbf{Cắt khoảng trắng thừa (Trim Whitespace)}: Loại bỏ các khoảng trắng ở đầu/cuối chuỗi.

\textbf{Chuẩn hóa chữ hoa/chữ thường (Standardize Case)}: Chuyển văn bản sang chữ thường hoặc chữ hoa để bảo đảm tính nhất quán.

\textbf{Sửa lỗi chính tả (Fix Spelling Errors)}: Các gợi ý AI của DataPrep có thể tự động sửa những lỗi gõ phổ biến (ví dụ: ``Smrtphone'' \ensuremath{\rightarrow} ``Smartphone'').

\textbf{Tách hoặc gộp cột (Split or Merge Columns)} (nếu cần): Ví dụ, nếu thành phố và bang được gộp trong cùng một trường, bạn có thể tách chúng thành hai cột.

Xem trước thay đổi

DataPrep hiển thị \textbf{ảnh chụp trước và sau} của dữ liệu mỗi khi bạn áp dụng một phép biến đổi. Hãy xác nhận kết quả đúng trước khi hoàn tất.

3. Nhóm giá trị thông minh cho các trường phân loại

Xử lý các nhãn không nhất quán

Với các danh mục như ``Product\_Category'', DataPrep có thể nhóm các giá trị đồng nghĩa hoặc gần như trùng lặp.

Nhấp vào tiêu đề cột và chọn \textbf{``Group Similar Values'' (Nhóm các giá trị tương tự).}

Phê duyệt hoặc tinh chỉnh các nhóm

DataPrep sẽ đề xuất gộp các giá trị như ``Smart Phone'', ``smartphone'' và ``SmarPhone''.

Hãy điều chỉnh thủ công hoặc từ chối các đề xuất nếu chúng gộp nhầm những giá trị khác biệt (ví dụ: ``Smartphone'' so với ``Smartwatch'').

Tạo tên danh mục nhất quán

Sau khi được phê duyệt, công cụ áp dụng các nhãn đã cập nhật cho toàn bộ cột, bảo đảm các danh mục đồng nhất trên toàn bộ tập dữ liệu.

4. Kiểm tra chất lượng dữ liệu tự động

Chạy quét ``Data Health'' (Sức khỏe dữ liệu)

Bản quét này phát hiện các bất thường như điểm đánh giá nằm ngoài khoảng cho phép (ví dụ: ``12'' trong thang đánh giá 1--10) hoặc định dạng email không hợp lệ.

Bạn sẽ thấy các vấn đề được gắn cờ trong thanh bên hoặc bảng tóm tắt lỗi.

Xem xét và áp dụng các đề xuất sửa lỗi

\textbf{Chuẩn hóa điểm đánh giá}: Nếu một điểm đánh giá là ``12'', bạn có thể chọn đặt thành ``10'' hoặc đánh dấu là không hợp lệ.

\textbf{Sửa định dạng ngày}: DataPrep có thể phân tích các định dạng ngày hỗn hợp (ví dụ: ``MM/DD/YYYY'' so với ``DD/MM/YYYY'') và chuyển chúng về một mẫu chuẩn hóa duy nhất.

\textbf{Sửa email không hợp lệ}: Công cụ có thể loại bỏ hoặc sửa các lỗi chính tả tên miền như ``@gamil.com''.

Đánh dấu các vấn đề dai dẳng là ngoại lệ

Nếu một số mục bị gắn cờ thực chất là các trường hợp biên hợp lệ, hãy đánh dấu chúng là chấp nhận được thay vì là lỗi.

5. Làm giàu và xác thực dữ liệu

Cột dẫn xuất

Tạo các cột mới từ các trường hiện có (ví dụ: tính ``Nhóm tuổi'' dựa trên ngày sinh).

Thêm quy tắc để xử lý các trường hợp ranh giới (ví dụ: nếu ngày sinh trước năm 1900, đánh dấu là ``Lỗi tiềm ẩn'').

Xác thực mã bưu chính

Với dữ liệu nhân khẩu học, hãy bật các kiểm tra theo khu vực (ví dụ: mã ZIP của Hoa Kỳ so với mã bưu chính của Canada).

DataPrep có thể làm nổi bật các điểm không khớp, nhắc bạn sửa hoặc loại bỏ các bản ghi không hợp lệ.

Tính nhất quán về thời gian

Chuẩn hóa các cột ngày/giờ về ISO 8601 (YYYY-MM-DD HH:MM:SS) hoặc một định dạng nhất quán khác.

Nếu tập dữ liệu của bạn có nhiều cột dấu thời gian (ví dụ: ``Response\_Submit\_Time'' và ``Survey\_Start\_Time''), hãy bảo đảm chúng dùng cùng một độ lệch múi giờ.

6. Tạo công thức (recipe) tái sử dụng và xuất dữ liệu

Lưu các phép biến đổi dưới dạng công thức (recipe)

Mỗi bước, tách cột, nhóm giá trị hoặc sửa kiểu dữ liệu, đều được tự động ghi lại trong một ``recipe''.

Đặt tên và đánh số phiên bản cho recipe của bạn để tái sử dụng hoặc cộng tác trong tương lai.

Xem trước tập dữ liệu đã làm sạch

Dùng chế độ xem mẫu để xác nhận mọi phép biến đổi đều đúng.

Kiểm tra số lượng bản ghi, tên cột và mọi trường phái sinh.

Tùy chọn xuất dữ liệu

Xuất dữ liệu đã làm sạch dưới dạng \textbf{CSV}, \textbf{Excel}, hoặc gửi trực tiếp đến một \textbf{cơ sở dữ liệu} hoặc bảng \textbf{BigQuery}.

Giữ lại recipe để các lần chạy sau có thể áp dụng cho các tập dữ liệu được cập nhật với rất ít công sức bổ sung.

Vì sao nên dùng Google DataPrep?

\textbf{Tích hợp với Google Cloud}: Kết nối dễ dàng với Cloud Storage hoặc BigQuery để xử lý dữ liệu có khả năng mở rộng.

\textbf{Gợi ý do AI hỗ trợ}: Các phép biến đổi thông minh giảm bớt phỏng đoán, giúp bạn tập trung vào các quyết định đặc thù theo lĩnh vực.

\textbf{Dấu vết kiểm toán \& Kiểm soát phiên bản}: Mọi phép biến đổi đều được ghi lại, giúp tăng khả năng tái lập và hỗ trợ phản biện đồng cấp (peer review).

\textbf{Xử lý tập dữ liệu lớn}: DataPrep có thể xử lý hàng triệu dòng mà không phụ thuộc vào bộ nhớ của máy cục bộ, điều này rất cần thiết cho các nghiên cứu quy mô doanh nghiệp hoặc nghiên cứu dọc.

\end{examplebox}

\section{Tích hợp và hài hòa hóa dữ liệu}

Trong nghiên cứu đương đại, dữ liệu hiếm khi được lấy từ một kho lưu trữ duy nhất hay được tạo ra bởi một công cụ đơn lẻ. Thay vào đó, chúng thường xuất phát từ các cơ sở dữ liệu, API, báo cáo lưu trữ và nguồn cấp dữ liệu thời gian thực rời rạc, mỗi nguồn có cấu trúc và ngữ nghĩa riêng (Kitchin, 2014). Việc tích hợp các nguồn không đồng nhất này thành một tập dữ liệu thống nhất, để các phân tích có thể nắm bắt toàn bộ phạm vi của một hiện tượng, đòi hỏi một quy trình cẩn trọng gồm đối sánh, ánh xạ và hài hòa hóa các trường dữ liệu rời rạc (Verborgh et al., 2016). AI, kết hợp với các nền tảng thân thiện với người dùng, có thể đơn giản hóa nhiệm vụ này bằng cách tự động phát hiện mối quan hệ giữa các tập dữ liệu và gợi ý cách tốt nhất để hợp nhất, căn chỉnh hoặc dung hòa các thành phần lược đồ (schema) xung đột.

Trong khi kỹ năng lập trình và viết mã lệnh từ lâu vẫn là rào cản đối với việc tích hợp dữ liệu phức tạp, ngày càng nhiều công cụ \textbf{no-code} hoặc \textbf{low-code} đang giúp những khả năng này dễ tiếp cận hơn (Ciechanowski, Jemielniak, and Gloor 2020). Các nền tảng như \textbf{Zapier}, \textbf{Make (trước đây là Integromat)} và \textbf{Airtable} cho phép nhà nghiên cứu kết nối các dịch vụ khác nhau, như Google Sheets, cơ sở dữ liệu SQL và thậm chí cả các API web, thông qua giao diện kéo và thả trực quan. Các công cụ này thường tích hợp những đề xuất do AI điều khiển để ánh xạ và hợp nhất các trường dữ liệu; chẳng hạn, chúng có thể gợi ý nối hai bảng dựa trên các tên cột tương tự hoặc nhắc người dùng căn chỉnh định dạng ngày tháng giữa các nguồn dữ liệu. Khi được sử dụng cẩn trọng, những gợi ý tự động này có thể giảm đáng kể thời gian tích hợp.

Khi việc tích hợp và hài hòa hóa dữ liệu trở nên đơn giản hơn nhờ các nền tảng no-code hoặc low-code, nhà nghiên cứu có thể dành nhiều sự chú ý hơn cho việc diễn giải kết quả và thiết kế các phân tích tiếp theo, thay vì sa lầy trong các thao tác hợp nhất và biến đổi thủ công. Tuy nhiên, việc xác thực các gợi ý tự động, duy trì kiểm soát phiên bản của các tập dữ liệu đã hợp nhất và lưu giữ tài liệu chi tiết để bảo đảm tính minh bạch vẫn là điều thiết yếu (Kaisler et al., 2013). Khi được triển khai một cách chu đáo, các quy trình tích hợp do AI hỗ trợ giúp giảm sai sót, rút ngắn thời gian từ dữ liệu đến hiểu biết sâu sắc (time-to-insight) và trao quyền cho các nhóm đa ngành kết hợp và dung hòa các nguồn dữ liệu đa dạng hiệu quả hơn.

\begin{examplebox}{Ví dụ thực tiễn}

Hãy tưởng tượng bạn có một khảo sát trực tuyến trên \textbf{Typeform} thu thập các câu trả lời dạng văn bản tự do, cùng một nguồn thứ hai là tệp CSV xuất từ hồ sơ nhân khẩu học của cơ sở bạn. Một nền tảng như \textbf{Zapier} có thể tiếp nhận các phản hồi khảo sát mới, đối chiếu các địa chỉ email trùng khớp trong tệp CSV và hợp nhất dữ liệu thành một bảng duy nhất trong \textbf{Airtable}. Các nhà nghiên cứu còn có thể tận dụng các gợi ý do AI hỗ trợ trong tính năng ``Automations'' của Airtable để thống nhất các nhãn nhân khẩu học không nhất quán hoặc sửa các địa chỉ bị cắt cụt trước khi lưu các bản ghi đã cập nhật.

Nếu bạn thường xuyên nhận nhiều tệp Excel từ các cộng tác viên khác nhau, bạn có thể dùng \textbf{Make} để theo dõi một thư mục chia sẻ trên Google Drive, tiếp nhận từng tệp được gửi đến và nối các hàng của tệp đó vào một tập dữ liệu chủ (ví dụ trong Google BigQuery). Hệ thống sẽ gắn cờ mọi cột mới hoặc kiểu giá trị bất thường, để bạn quyết định cách tích hợp chúng.

Đôi khi, dự án của bạn bao gồm các bản ghi chép phỏng vấn dạng văn bản cùng dữ liệu khảo sát dạng số. Các công cụ như \textbf{Power Automate} (thuộc Microsoft Power Platform) có thể điều phối các luồng dữ liệu, trong đó một bước phân tích văn bản do AI hỗ trợ (ví dụ Azure Cognitive Services để nhận dạng thực thể) được chạy trên các bản ghi chép. Các kết quả đầu ra đã xử lý, như điểm cảm xúc hoặc nhãn chủ đề, được nối vào một tập dữ liệu trung tâm vốn cũng chứa các câu trả lời khảo sát của từng người tham gia. Cách tiếp cận này giúp đơn giản hóa các phân tích tương quan giữa các nhận xét mở và các đánh giá có cấu trúc.

Nếu bạn có nhiều tập dữ liệu gán nhãn khác nhau cho cùng một biến, chẳng hạn ``SatisfactionScore'', ``User\_Sat'' và ``Sat\_Score'', các gợi ý do AI hỗ trợ trong một công cụ như \textbf{Airtable} hoặc \textbf{OpenRefine} có thể đề xuất một tên trường thống nhất, đã được hài hòa hóa. Bạn chỉ cần chấp thuận đề xuất, và nền tảng sẽ áp dụng thay đổi cho toàn bộ bản ghi, giúp giảm đáng kể công sức xử lý thủ công.

Các nhà nghiên cứu cũng có thể khai thác các nguồn dữ liệu bên ngoài, chẳng hạn API của một nền tảng mạng xã hội hoặc nguồn cấp dữ liệu thời tiết. Khi dùng một công cụ tích hợp no-code, bạn cấu hình một điểm cuối (endpoint) để lấy dữ liệu JSON hoặc XML, rồi để các tính năng ánh xạ do AI hỗ trợ đề xuất cách căn chỉnh các trường đầu vào với các cột trong cơ sở dữ liệu hiện có của bạn. Hệ thống có thể nhận ra rằng ``temp\_f'' từ nguồn cấp thời tiết nên khớp với cột ``Temperature\_Fahrenheit'' của bạn, hoặc ``user\_screen\_name'' từ API của X ánh xạ tới ``TwitterHandle''. Việc xác nhận hoặc điều chỉnh các ánh xạ này bảo đảm tên gọi và kiểu dữ liệu nhất quán trên toàn bộ tập dữ liệu tích hợp của bạn.

\end{examplebox}

\section{Phân tích khám phá và kỹ thuật đặc trưng với AI}

Phân tích dữ liệu khám phá (Exploratory Data Analysis, EDA) nhằm làm lộ ra các mẫu hình tiềm ẩn, các điểm bất thường và các mối quan hệ trước khi quá trình xây dựng mô hình chính thức bắt đầu (Tukey, 1977). Theo truyền thống, các nhà nghiên cứu phải tự tay lọc qua các biểu đồ và thống kê tóm tắt, nhưng những tập dữ liệu đồ sộ ngày nay khiến công việc này trở nên đáng ngại (Han, Kamber, and Pei, 2012). Các công cụ dựa trên AI có thể tự động hóa phần lớn công việc nặng nhọc, gợi ý những mối tương quan quan trọng và làm nổi bật các phân phối đáng chú ý. Chẳng hạn, hãy hình dung một tập dữ liệu gồm 10.000 hồ sơ học sinh từ nhiều trường khác nhau trên khắp cả nước: một nền tảng no-code như \textbf{Tableau Prep} có thể nạp tệp CSV của bạn, tự động tạo biểu đồ histogram cho từng biến và làm nổi bật các giá trị ngoại lai trong tỷ lệ đi học. Đồng thời, một quy trình làm việc với \textbf{ChatGPT} có thể tạo ra những bản tóm tắt ở mức tổng quát (``Tỷ lệ đi học có tương quan cao với điểm thi cuối kỳ, đặc biệt ở các khu vực có thu nhập thấp hơn''), giúp bạn khỏi phải rà soát thủ công hàng trăm bảng tổng hợp (pivot table).

Một lợi ích của các nền tảng EDA dựa trên AI này là chúng có thể xử lý các tác vụ như phân tích tương quan, phân cụm và giảm chiều dữ liệu mà hầu như không cần lập trình (Witten, Frank, and Hall, 2011). Chẳng hạn, một nhà nghiên cứu dữ liệu khí hậu có thể tải các tệp chuỗi thời gian lên giao diện no-code của \textbf{DataRobot}, từ đó tự động tạo ra các bản đồ nhiệt tương quan và biểu đồ phân tán cho những biến như nhiệt độ, độ ẩm và mức độ ô nhiễm. Chỉ với vài cú nhấp chuột, bạn có thể thấy độ ẩm tăng vọt đồng thời với một số chỉ báo ô nhiễm cụ thể, những hiểu biết có thể không hiển nhiên cho đến khi một quy trình tự động chỉ ra chúng. Sau đó, bạn có thể tinh chỉnh những phát hiện này bằng các truy vấn nâng cao trong một môi trường thân thiện với người dùng, thay vì viết mã mọi thứ từ đầu.

Chuyển sang kỹ thuật đặc trưng, các nhà nghiên cứu thường xuyên tạo ra các biến mới, bằng cách kết hợp các cột hiện có hoặc trích xuất những tín hiệu có ý nghĩa từ dữ liệu thô, nhằm nâng cao hiệu suất và khả năng diễn giải của mô hình (Kuhn and Johnson, 2013). Trong một tình huống marketing, bạn có thể có nhật ký mua hàng của khách hàng chứa các dấu thời gian. Sử dụng một công cụ trực quan như \textbf{Airtable} với các tiện ích mở rộng tích hợp AI, bạn có thể tạo cột ``TimeSinceLastPurchase'' hoặc suy ra đặc trưng mua hàng ``Buổi sáng so với Buổi tối'' chỉ bằng cách nhấp qua một quy trình có hướng dẫn. Nếu bạn cần các phép biến đổi văn bản, ví dụ, chuyển các đánh giá sản phẩm viết tự do thành điểm cảm xúc, một prompt (câu lệnh) nhanh gửi tới \textbf{ChatGPT} có thể tạo ra các ví dụ regex hoặc mã giả mà bạn có thể điều chỉnh để dùng trực tiếp trong môi trường no-code của mình.

Kỹ thuật đặc trưng tự động có thể đi xa hơn bằng cách đề xuất theo thuật toán các phép biến đổi, phân nhóm (binning) các giá trị số, hoặc tạo ra các biến tương tác (Chandrashekar and Sahin, 2014). Giả sử bạn đang nghiên cứu một tập dữ liệu y tế công cộng với số bước chân hằng ngày, chỉ số khối cơ thể (BMI) và số giờ ngủ. Một nền tảng như \textbf{H2O Driverless AI} có thể đề xuất các phép biến đổi, chẳng hạn ``số bước chân trung bình trong 7 ngày qua'' hoặc ``nhóm BMI nhân với số giờ ngủ trung bình'', có thể nâng cao khả năng dự báo đối với một số kết cục sức khỏe cụ thể. Ngay cả khi bạn không tự tin về thống kê nâng cao, những khuyến nghị này cho phép bạn thử nghiệm nhanh với các biến mới và nhận phản hồi tức thì về việc chúng ảnh hưởng thế nào đến độ chính xác hoặc khả năng giải thích của mô hình.

Một trường hợp sử dụng thực tế khác liên quan đến việc làm việc đồng thời với dữ liệu văn bản và dữ liệu số. Hãy hình dung bạn đã thu thập các câu trả lời khảo sát mở về mức độ hài lòng với sản phẩm cùng với điểm đánh giá theo số sao. Một công cụ như \textbf{MonkeyLearn} có thể tạo điểm phân tích cảm xúc cho từng bình luận, trong khi bảng điều khiển no-code của bạn, chẳng hạn \textbf{Power BI}, tự động gộp các điểm này với các điểm đánh giá bằng số trong một bảng duy nhất để phân tích tương quan khám phá. Bạn có thể phát hiện rằng các đánh giá ít sao thường mâu thuẫn với những cảm xúc trong văn bản lại tích cực một cách đáng ngạc nhiên, từ đó thúc đẩy việc điều tra sâu hơn về thiết kế khảo sát hoặc trải nghiệm người dùng. Sự cộng hưởng giữa EDA dựa trên AI và kỹ thuật đặc trưng đảm bảo rằng ngay cả những người không biết lập trình cũng có thể nhanh chóng khám phá những hiểu biết đa chiều (Witten, Frank, and Hall, 2011).

Cần nói thêm rằng Google Sheets mặc định cung cấp nhiều công cụ làm sạch dữ liệu, và có rất nhiều tiện ích bổ sung thực hiện các phân tích bổ sung, bao gồm cả phân tích cảm xúc. Tuy nhiên, thường rất khó để kéo dữ liệu từ một tập dữ liệu lớn hơn. Một tiện ích bổ sung hữu ích và đáng chú ý cho Google Sheets là Wikipedia and Wikidata Tools, cho phép truy vấn các cơ sở dữ liệu của Wikimedia ngay từ bên trong bảng tính.

Bản thân việc sử dụng dữ liệu Wikimedia cũng vô cùng hữu ích cho phân tích dữ liệu khám phá. Các đồ thị tri thức tự do và mở quy mô lớn, được duy trì cộng tác, như \textbf{Wikidata}, cho phép thực hiện những nghiên cứu và trực quan hóa tuyệt vời, nhanh chóng mà không hề sơ sài (Turki et al.~2022, 2021). Wikidata chứa dữ liệu có cấu trúc, được liên kết về hàng triệu thực thể, con người, địa điểm, khái niệm, v.v., khiến nó trở nên lý tưởng để bổ sung ngữ cảnh hoặc các mối quan hệ cho một tập dữ liệu hiện có. Chẳng hạn, một nhà nghiên cứu khảo sát sự phân bố địa lý của những người đoạt giải Nobel có thể truy vấn Wikidata để lấy siêu dữ liệu chuẩn hóa về quốc gia sinh, đơn vị công tác học thuật hoặc lĩnh vực nghiên cứu của từng laureate. Sau đó, họ có thể gộp các chi tiết này vào một nền tảng phân tích no-code như \textbf{Tableau} hoặc \textbf{Power BI} để trực quan hóa thêm. Wikidata và Wikipedia cũng có thể cho phép thực hiện những nghiên cứu thú vị: một người trong chúng tôi đã phát hiện rằng các vận động viên đoạt huy chương Olympic có tuổi thọ ngắn hơn so với các vận động viên Olympic không đoạt huy chương (Kovbasiuk, Ciechanowski, and Jemielniak, 2024) bằng cách khai thác các nguồn tài nguyên miễn phí của Wikimedia.

Các công cụ như \textbf{Wikidata Query Service} (sử dụng SPARQL) cung cấp giao diện thân thiện với người dùng để truy xuất những lát cắt cụ thể của dữ liệu liên kết. Đồng thời, các trình bao bọc (wrapper) Python (ví dụ: \textbf{wikidata2df}) có thể tự động hóa các truy vấn này cho những quy trình làm việc có thể lặp lại. Việc tích hợp liền mạch các nguồn tri thức mở này có thể làm lộ ra những mối liên hệ ẩn, chẳng hạn các trường đại học chung hoặc những chủ đề nghiên cứu chồng lấn, mà nếu không thì có thể bị bỏ sót trong các tập dữ liệu cục bộ, qua đó làm tăng chiều sâu và độ rộng của phân tích khám phá.

Bất chấp những ưu điểm này, sự thận trọng là điều thiết yếu. Các quy trình tự động có thể đề xuất một số lượng lớn biến mới, trong đó một số biến có thể phản ánh những mô thức nhất thời hoặc giả tạo thay vì các tín hiệu thực sự (Friedman, 2001). Khi bạn phát hiện một mối liên hệ bất ngờ, chẳng hạn mối tương quan mạnh một cách kỳ lạ giữa ``Average Customer Ticket Size'' và ``Browser Language'', thì nên kiểm chứng nó dựa trên tri thức chuyên ngành và có thể kiểm định chéo kết quả bằng một cuộc kiểm tra thủ công chặt chẽ hơn. Bằng cách kết hợp hiệu quả của EDA và kỹ thuật đặc trưng do AI điều khiển với phán đoán học thuật có hiểu biết, các nhà nghiên cứu có thể đẩy nhanh quá trình khám phá đồng thời giảm rủi ro chạy theo nhiễu. Về bản chất, các công cụ AI không thay thế sự thấu hiểu của con người; chúng trao thêm sức mạnh cho sự thấu hiểu ấy, cho phép các nhà nghiên cứu dành nhiều năng lượng hơn cho những đột phá khái niệm và ít hơn cho các tác vụ dữ liệu lặp đi lặp lại.

\section{Kết luận}

AI đang làm thay đổi căn bản việc quản lý và phân tích dữ liệu, giúp các nhà nghiên cứu xử lý hiệu quả những bộ dữ liệu ngày càng lớn, phức tạp và không đồng nhất. Bằng cách tự động hóa các tác vụ tẻ nhạt nhưng thiết yếu, như xử lý và chuẩn hóa dữ liệu (data wrangling), làm sạch, tích hợp và phân tích khám phá, các công cụ dựa trên AI nâng cao đáng kể năng suất, độ chính xác và khả năng tái lập của nghiên cứu. Những công nghệ này giúp học giả nhanh chóng phát hiện các mô thức ẩn, sửa lỗi, hài hòa các nguồn dữ liệu rời rạc và xây dựng những đặc trưng mới giàu thông tin, qua đó tạo điều kiện cho các phân tích sâu sắc và tinh tế hơn. Tuy nhiên, triển vọng của việc quản lý dữ liệu dựa trên AI luôn phải được cân bằng bằng sự kiểm chứng nghiêm ngặt, các cân nhắc đạo đức và việc ghi chép minh bạch. Dù AI tinh gọn quy trình làm việc và nhận diện những hiểu biết có giá trị, phán đoán của con người vẫn giữ vai trò then chốt trong việc xác minh các gợi ý tự động, ngăn chặn sự lan truyền của các thiên lệch mang tính hệ thống và bảo đảm tính giá trị của các kết luận khoa học. Rốt cuộc, AI không thay thế chuyên môn học thuật mà bổ trợ đáng kể cho chuyên môn đó, giải phóng nhà nghiên cứu khỏi các tác vụ nhàm chán để họ gắn bó sâu hơn với dữ liệu của mình, thúc đẩy những khám phá mang tính đổi mới và nâng cao tính chặt chẽ chung của nghiên cứu học thuật.
